\section{Geometric Foundations and Continuous Trajectory Theory}
\label{sec:theory}

\subsection{Sequence Metric Space and Buneman Coordinates}
\label{subsec:metric_space}

We ground the method in an authentic viral benchmark: the HIV-1 circulating recombinant forms CRF07\_BC and CRF08\_BC \cite{su2000characterization}.
Both lineages emerged from independent recombination events between Subtype B and Subtype C in China.
They share parental sources across the 3,120-nucleotide \textit{pol} gene.
Real chimeras test theory.
Our alignment comprises 22 sequences: eight Subtype B references, eight Subtype C references, three CRF07\_BC isolates, and three CRF08\_BC isolates.
From these sequences, we compute the global pairwise evolutionary distance matrix $\mathbf{D} \in \mathbb{R}^{22 \times 22}$ under the Tamura-Nei substitution model with pairwise deletion of gap characters, capturing nucleotide transition and transversion rate biases across the retroviral genome \cite{tamura1993estimation}.
We center the squared distances using the projection operator $\mathbf{H} = \mathbf{I}_N - \frac{1}{N} \mathbf{1}_N \mathbf{1}_N^T$, forming the inner-product Gram matrix $\mathbf{B} = -\frac{1}{2} \mathbf{H} \mathbf{D}^{(2)} \mathbf{H}$ \cite{gower1966some}.
Spectral decomposition of $\mathbf{B}$ yields non-negative eigenvalues $\mathbf{\Lambda}_k$ and eigenvectors $\mathbf{V}_k$, projecting each sequence $i$ into a continuous coordinate $\mathbf{x}_i = \mathbf{V}_{k, i} \mathbf{\Lambda}_k^{1/2} \in \mathbb{R}^d$ \cite{buneman1974note}.
The primary axis cleanly separates Subtype B from Subtype C.
Canonical metric coordinates are computed once for the entire alignment.
An invariant coordinate basis eliminates Procrustes alignment overhead and establishes a fixed reference frame for continuous trajectory tracking.

To track sequence affinity along the chromosome, each genome requires an instantaneous coordinate at every position $u \in [1, L]$.
Recalculating multidimensional scaling across sliding windows would distort the geometric space and require expensive Procrustes rotations.
Instead, we project local distance vectors into the fixed coordinate basis using Gower's analytical out-of-sample mapping \cite{gower1968adding}.
At coordinate $u$, sequence $i$ presents an instantaneous distance vector $\mathbf{d}_i(u) \in \mathbb{R}^N$ to the $N$ reference genomes, projecting to $\mathbf{z}_i(u) = \frac{1}{2} \mathbf{\Lambda}_k^{-1/2} \mathbf{V}_k^T (\bar{\mathbf{d}}^2 - \mathbf{d}_i^2(u))$ where $\bar{\mathbf{d}}^2 = \frac{1}{N} \mathbf{D}^{(2)} \mathbf{1}_N$ represents the mean squared reference distance.
When applied to our HIV-1 benchmark, the resulting coordinate trajectories reveal the physical dynamics of genetic exchange.
Non-recombinant Subtype B and Subtype C genomes remain stationary within their respective corridors across the entire gene.
Wobble never crosses clades.
Point mutations generate harmless local flutter around parental medoids, but unreticulated lineages never breach the inter-subtype divide.
In sharp contrast, the recombinant CRF07\_BC trajectory tracks within Subtype C territory through Protease, abruptly vaults across the zero-axis into Subtype B space in Reverse Transcriptase, and plunges back into Subtype C.
CRF08\_BC shares the identical parental gene pool but enters the Subtype B cassette approximately 130 nucleotides upstream.
In two-dimensional metric space, the two chimeras trace distinct orbital loops, resolving competing recombinant structures in linear time without triplet enumeration.
Manifolds preserve chimeras.
